\chapter{From \$5 Cloud Node to Bastion Host}


Cloud providers hook us in with cheap teaser rates; and then they slowly offer us feature after feature, at \$10 each.  Before you know it, the \$5 a month budget has grown to \$50, \$60, or \$70.  We'll tell the story of how we can took a bare bones FreeBSD VM from a major cloud hosting company and built it up to be a standalone bastion host that works as a nameserver and more.  


\section{Our Goals}
For this project, we wanted to buy a domain name, and then do everything else ourselves.  Aside from needing the registrar to get the name, we didn't want to have to use any of the provided services.  Registrars can do a great job of offering web hosting, site building, and other services; but, we shouldn't have to depend upon them.  If we took the lowest cost deal they offered for hosting, what could we do?


For a projected cost of \$10 a month, we could get two VMs.  These nodes would each have their own public IP address.  They'd be bare bones FreeBSD systems.  Being "out there on the Internet" would be to our advantage; sites acting as nameservers would serve us better away from our main node.  Since our main demarc only gets one IP, these two VMs might be a good fit.

\section{Domain Name Purchasing}
We bought a pair of names that we like.  We didn't buy any extras.  We didn't get hosting, \gls{certificate}s, email:  none of that.  Just get a name.


\section{Before We Built Anything We Decided on Keys}

With an account created at the hosting provider, we could quickly see that they offered us just enough power to get ourselves in trouble.  First off, they wanted us to set either a root \gls{password} or upload an SSH key.  Choose the key.  Building an SSH key is a little more involved; but, we don't want to arrive at the node, in the root account via SSH, with a simplistic \gls{password}.

That's right:  the hosting provider, by default, chose to provide SSH into root, first thing.  Now, after a while, we'd all do what we could to plus up the node and improve things.  However, the Internet is the Wild West.  Already out there will be scanning nodes looking for new accounts like ours.  It'd be best to arrive tough enough to repel common attacks.  For this reason, we recommend setting up an SSH key if you have the choice.

Since we chose the SSH key, we were not able to use the provided web console to look into our host.  You guessed it:  if you pick a public key, then you have to use it.  So, we did.  

\section{Establishing a \index{BIND}BIND \index{DNS}DNS Server with MX to an Encrypted Email Provider}


We quickly established a \index{BIND}BIND9 \index{DNS}DNS server, using this tutorial: \endnote{[1]}

Our domain name provider had a set of dialogs that allowed us, through our account, to direct our domain name to our own nameserver.  Armed with the IPs from the VMs and some progress through the \index{DNS}DNS server tutorial, we were able to adapt and fill out the records wtih the registrar.

We stopped right before the "\index{DNSSEC}DNSSEC" portion of the tutorial and added our MX records.  For our email server, we chose a free online email provider that offers storage in Switzerland.  As part of our subscription with them, we were able to use custom domains.  To get that to work, we had to follow a set of dialogs the email provider had for setting up critical \index{DNS}DNS records for SPF, DKIM, and DMARC.  It took a little bit for our changes to propagate, but it worked with little fuss.


Once our nameserver was found and our emails were arriving at the desired account, we walked through \index{DNSSEC}DNSSEC.  This was one advantage we had over what the registrar offered:  they don't do \index{DNSSEC}DNSSEC setups.  With the tutorial, we got it working in less than an hour.  By carefully following the directions, we were able to get \index{DNSSEC}DNSSEC working right away.  It was actually one of the easiest aspects of the system deployment.

\section{Dynamic DNS and Our Node}
With two nodes out on the Internet, and one nearby, there was always the chance that a service provider might change our address.  In practice, we were able to get stable addresses just by asking our local muni fiber company to provide one.  Still, they could change at any time.  So, we used an existing subscription to a dynamic DNS provider to address the hosts.  

To use a dynamic DNS with FreeBSD, we install ddclient.  Our service provider gave us example \gls{configuration} files that worked on the first try.  After a while, we began to realize that, if our nameserver was working well, then there was no reason why we couldn't run a dynamic DNS service of our own.  ddclient will work with RFC compliant procedures.  We found tutorials to support that innovation on FreeBSD, too.


It was just easier to have some addresses with simple names, through dynamic DNS.  Time and again, during configurations like these, we would need to call up the computer.  From one to the other and back again, we'd type in calls.  It can be helpful and convenient during the confusing moments of \gls{configuration} to be able to call in to your desired host easily.  

In the case of whitelisting hosts in \gls{firewall}\index{firewall} rules, it can be helpful to also list some hosts using their dynamic dns domain name.  If the IPs change, and \gls{firewall} rules hold hardcoded IP addresses, then how will you get back in to the host you've just secured?  Dynamic \index{DNS}DNS addresses in the \gls{firewall} rules offer a saftey during the confusing moments of \gls{configuration}.

\section{Firewall Rules with pf}


Pf comes installed as part of the base \gls{operatingsystem} in FreeBSD.  There are also two others we could choose from.  It's mostly a matter of preference.  We picked pf.  

To set up the rules, we consulted some tutorials.  Early on, we wrote pass in rules for our favorite IPs and those dynamic DNS host and domain names, as mentioned above.  

We soon found out we were locking out critical services.  By using searches of /etc/service for keywords, we often found \gls{port}s that needed to be kept open to keep our machine running smoothely.  Since we were drafting some of these rules for the first time, this took some expriementation.  

When setting up \gls{firewall} rules and login limitations, it's helpful to use more than one \index{ssh}ssh session.  Use one session to maintain a connection and adjust the config.  Use another to test access to the host.  If you don't, and you mess up, then you might find yourself locked out.  

\section{External Monitoring with Shodan}


Shodan.io offers a monitoring service.  It's free with the developer's account, up to about a dozen hosts.  Around 14 or 15, they start to require a paid Enterprise account.  Since we have few computers to look after, we snapped up Shodan Monitor.  

With Shodan, we learned the painful history of the past users of our IP address.  We also learned about the locations and open \gls{port}s of whomever it was who was calling us on \gls{port} 18888 over 2,500 times in 20 minutes.  One IP had once been misused as an Internet scanner.  The other seemed to be getting scanned constantly.  It's a good thing we activated pf.  Until we did, we didn't have easy visibility of who was calling the node unsuccessfully.  

From the inside, most hosts look quiet.  Once we set up a \gls{firewall}, we can see who tries to \index{SSH}SSH in.  We can see who sends hacky attempts to rattle our server.  And, with Shodan, we can get an idea about what computer their using to do it.  "IP is not ID," but it's always interesting to see what machine is calling us.


\section{Proxy:  Providing Standoff Distance and a Planned Access Approach}
With the DNS server up and running, we had a lot of space left over. The minimal, bare-bones host provided by the cloud service had used barely 10% of its capacity.  We wondered what else we could do with the machine.  We could screen.

With a \gls{firewall} already running (and repelling unwanted calls) on the master and slave \index{DNS}DNS servers, wouldn't it be nice to use them as standoff entrances for our site?  With a demarc to safeguard, it would be best if we could require traffic to enter at the distant hosts and then walk down to our site along a path we could specify.  This seemed to be the best way to build in "standoff distance" into the acces plan.  

With standoff distance, we have a link between our data center demarc and the general public that we can monitor, secure, and modify.  What if we have a DoS attack?  If our front door is our only door to the Internet, we won't have any room to move or leverage to create a response.  Again, having a VM on the Internet in a distant place is to our advantage.  

Now, as traffic arrives at our data center demarc, it'll hit a \gls{firewall} right away.  Later, as it moves down the line to our desired server, it'll hit some auth.  So, inside our data center we could react to traffic; but, the proxies on the distant hosts will help us control traffic that arrives to the server that's our goal.  At the data center demarc, we can have a \gls{firewall} rule that will refuse traffic for our desired server unless it walks down the path from our proxy.

\section{Installing HAProxy}


Installing the program was a snap.  We had to read some documentation, and follow some suggestions; but, it all worked without a hitch.  There are several techniques we can choose from with HAProxy.  Ultimately, we chose 
TCP Forwarding.  This is a "layer 4" proxy.  HAProxy has OSI Layer 7 capabilities that can let us filter traffic based on contents.  With tcp mode, haproxy will just pass the traffic on down to our data center demarc, invisibly.

SSL stayed on the web host.  Before we began with these \index{DNS}DNS servers, we had already installed Let's Encrypt TLS \gls{certificate}s on our target webserver.  By using tcp mode, we could still serve up our pages without interfering with the TLS integrity.  Some other proxy techniques might have required TLS termination.  This would have meant re-installing certs on the proxy hosts in order to present a clean connection to the user.  Tcp mode was simple, effective, and easy to use.

To check out our proxy install, we started with calling the nameserver host address.  It would then use haproxy to walk us down the line, reques our page from the webserver, and then the page we called would come back up the line.  Satisfied that the proxies were working, it was time to change the DNS records.

\section{DNS Load Balancing}


Most DNS records get served round-robin, by default.  That is, when multiple aliases (A records) are listed, they'll be shown in a rotating order by the nameserver who answers.  Since our proxies were now on the nameserver hosts, we could direct our traffic to each.  We changed our A Records for the domain to point to each DNS server, instead of our demarc IP.  This is the outside half of getting our traffic to walk down the funnel into our screen.  

As mentioned earlier, there are \gls{firewall} rules near the demarc to do the initial onsite screen to require incoming traffic to only be considered if it used the proxy.  

Notice how the DNS load balancing could have some unenforceable conditions.  For example, what if a received DNS record was cached?  Just because a URL is called many times by a user doesn't mean that it will get its values from our authoritative nameserver.  Any DNS server along the way might hold the answer.  After all, that 48 hour propagation, coupled with sometimes deceptive "instantaneous" DNS response from some servers, can be deceptive.  The propagation of DNS records does take time.  Just as wrong answers can be out there for a while, beyond our control; so also we can have a printing of DNS aliases in an order that was previously served and incompletely rotated because of caching.  DNS load balancing is more of a suggestion than a rule because our overall control of the records is gone once they leave our computer.

HAProxy offers load balancing.  In our case, we were only going to one destination, so most of those features didn't yet need to be activated.  They might be more useful on a haproxy install inside the demarc.  It's true load balancing.  We can set rules that will filter traffic by subdomain to specific servers.  We can ensure high availability with health checks on servers that precede traffic routing.  

\section{TOTP with Google Authenticator}


Remember our login situation?  We chose to use SSH public keys on the server, which was a good idea; but, we wanted to do a little more.  We added on Google Authenticator's time-based one time passwords (TOTP) to the ssh login accounts.  
FreeBSD ports have a package for this.  It's pam\_google\_authenticator.  Thanks to scant and incomplete documentation, setting up Google Authenticator on FreeBSD was a real bear.  When editing the sshd\_config, we have to choose carefully \gls{password}ssword options we want to keep and which we want to comment out.  On one hand, we're using ssh public key.  On the other, if we turn off all of the password\index{password} functions, TOTP won't work.  

When misconfigured, we can end up in a few different situations:
 - We never get a prompt for a verification code
 - We get repeatedly prompted for a code, but it never works
 - We get prompted for a public key, a verification code, and a \gls{password}.

This last one is a little bit of a problem because, as you might guess:  we didn't always set one.  Also, if you want to use Google Authenticator with a root account, there is one more setting in the config that needs to be hunted up and adjusted.  Most folks, unlike our hosting provider, don't want to ever let someone ssh into a root account.  Guess why.

Setting up the Google Authenticator on the host wasn't too hard, dialog-wise.  There is a short program that's run to ask some simple questions.  There is a 3D Q barcode that's shown.  It'll be printed in terminal with ASCII graphics.  The results of the dialog will be stored in a hidden file.  That includes the "scratch codes" to get in if the app fails.  This means that anyone with access to the hidden file might be able to view those codes.  Keep that in mind if you do a multi-user install.

Similar procedures can be used with other TOTP products and FreeBSD.  With our VM far away, we'll never be able to plug in a USB dongle.  Also, with programs like Google Authenticator, there are other programs that may not require the possession of a specific phone.  Keep these ideas in mind when assessing risk related to TOTP products.  Their fascinating to watch, but they have their limitations just like everything else.

\section{Backup and Restore with tar}


We've done enough work to not want to do it all again.  Time for a \gls{backup} and retore policy.  We'll need to make our own.

The hosting provider does weekly, automated backups.  They also offer snapshots of the host.  However, that doesn't quite meet our needs.  The hosting company retains control of and access to the backups.  We can apply them to the host through a web interface, but we can never download and store them on physical media we control.

They also offer snapshots.  Snapshots are diff files.  We can restore the VM to them; but, again they're not full backups that we can really control.  If things really go haywire, a diff file to another problem can be its own problem.  Snapshots are a convenient way to reset a VM, but they don't offer the comprehensive protection that's implied by the title, "Backup\index{backup}."  

Unlike dd a hard drive in our lab, we don't have physical access to these disks that hold the VM at the hosting company.  That means that we also can't swap them out, use them with another motherboard, or generally control the automated backups and snapshots.  We can rehearse a restoration on an offline machine in our lab with our backups; but we can't do that with just what the hosting company offers.  They're offsite compared to our place, and that's good; but, they're also not offline.  We can only access them through the online controls offered by the hosting provider.  Let's recognize this point:  the backups and snapshots offered through the web interface at the hosting provider don't meet the criteria of DFIR survivability we'd expect from a machine on the bench in our lab.  

This means that if we had to reconsitute or rebuild the host from scratch, we would be at a loss without the help of the hosting company.  Since this is not an acceptable situation, it's up to us to get a \gls{backup} done.  And, if we can remember, a \gls{backup} is not a backup until it's been restored.  


To give ourselves offsite, offline backups, we're going to use tar to take a sample of critical files used in the installation of the programs above.  One by one, we went through directories that held configs.  We built a script that would automatically copy them to a directory.  Then we tarred the achive, g-zipped it, and exfiltrated it with scp.  Once on a computer we could physically access, we saved the archive to removable storage.  Cataloged and moved to a safe place, our tar and restore sample was ready for a restoration rehearsal.

\section{Installing OSSEC-HIDS}

\begin{lstlisting}
You need to create main \gls{configuration} file:
/usr/local/ossec-hids/etc/ossec.conf

For information on proper configuration see:
https://www.ossec.net/docs/syntax/ossec_config.html
\end{lstlisting}

This was a little more tricky than usual.  Most of the tutorials and documentation out there right now encourages us to edit a file called ossec.conf.  However, there's been a change in the architecture of the system.  What was ossec.conf is now cobbled together from other, smaller files.  Those files start with names like "900.local.conf"; and, around those files, we'll see other smaller directories named after significant parts of ossec.conf.  In various files, there are warnings about not editing something.  It's very confusing.

There was no advice at the official ossec-hids sites about these changes.  Trying to figure this out was very painful.  Even the package installation notes themselves advised us to edit ossec.conf.  Trying to figure out how to adjust this thing was so confusing we almost abandoned the program.  Fortunately, we had installed one some time before; back then, the installation directions were clear and it was much easier.  We kept on and got it installed.

The end of it is:  place your edits in the 900.local.conf file.  It'll persist through reboots and upgrades that way.  That file will hold nested XML stating your preferences; it's just as ossec.conf would have had elements wedged into it if you had edited it the old way.

There are three kinds of installs for this program:  local, agent, and server.  Pick local for the first try.  This means that the rules and logs of monitoring will be installed on the same host that's getting monitored; but, you've got to start somewhere.  It's the easiest way to figure out the functions.  There are few good tutorials.  On an early Saturday morning, I found some documentation lacking.  It's a small project.  You'll need to figure out how to help yourself some.

In several of the tutorials I did see, it was frequently advised to download from pkg.  This presents another problem:  the default pkg install is for "no firewall\index{firewall}."  As you can see from the pf install, above, that's not going to work.  Rather than try to modify what's provided, it was easier to just build ossec-hids from ports on FreeBSD.  In the make configs that get presented in ncurses dialogs, you'll see the options for adding in pf \gls{firewall}.

\section{Adjusting sendmail for ossec-hids}
The primary way ossec-hids will communicate with your sysadmin from a local install will be through email.  This means that sendmail should be configured for sending mail out only.  Also notice that the zone files we previously set up don't have a good record for just this host that will be doing the transmitting.  In fact, it could come into conflict with those SPF, DKIM, DMARC records set up earlier.  Adjust your zone files accordingly.

The email function is not automatically ready for use.  It will need to be edited, because critical address values can't be known by the program before its installation.  

You will need to uncomment an entry in one file to get the mail:
/usr/local/ossec-hids/etc/ossec.conf.d/900.local.conf

Which holds the XML element that corresponds to the old rule for:
<global>
    <email\_notification>yes</email\_notification>
    <email\_to>sammy@example.com</email\_to>
    <smtp\_server>mail.example.com.</smtp\_server>
    <email\_from>sammy@example.com</email\_from>
</global>

\section{Adjusting pfSense for Incoming HAProxy Traffic}
Earlier, we set up haproxy.  With the zone files populated to direct traffic to the proxy, we adjusted firewall rules on the demarc to accept calls only from those proxies.  As we were bringing up the system, we needed to accept traffic that was sent directly to the demarc.  Once we were confident that DNS records had propogated to point to the bastion hosts, then we could restrict access through the demarc through those hosts with rules.  

In upcoming drafts, we may discuss:
- snort
- tcpdump for continuous packet capture for archiving traffic history.






\begin{lstlisting}
\end{lstlisting}
 \[[]\]
\subsection{Annotated Bibliography}


\endnote{[1]} \_\_\_\_\_.  INTERNET:   \url{https://blog.andreev.it/?p=4096}\endnote{\protect\url{https://blog.andreev.it/?p=4096}}

\markchapterendnotes
